\subsection{Ablaufbeschreibung eines HTTP/2-Clients über h2c}

Settings Austausch

\begin{enumerate}
    \item TCP-Handshake
    \subitem Client sendet SYN \textrightarrow{}Verbindungsanfrage (Pakete 3)
    \subitem Server antwortet mit SYN-ACK \textrightarrow{}Verbindungsbestätigung (Pakete 4)
    \subitem Client bestätigt mit ACK \textrightarrow{}Verbindung steht  (Pakete 5)
    \item HTTP/2 Connection Preface \textrightarrow{}Information über die HTTP/2-Unterstützung (Paket 6)
    \item SETTINGS Frame \textrightarrow{}Konfiguration der Verbindung (Paket 6)
    \item Window Update Frame \textrightarrow{}Erlaubt dem Server, mehr Daten zu senden (Paket 6)
    \item Server sendet SETTINGS Frame (Paket 10)
    \item Client bestätigt SETTINGS Frame (Paket 12)
    \item Server bestätigt SETTINGS Frame (Paket 13)
\end{enumerate}

HPACK-Kodierung

\begin{enumerate}
     \item Client sendet HTTP/2 HEADERS Frame mit der HPACK-kodierten Anfrage und nutzt dafür seine Kodierungstabelle
     \item Server empfängt HEADERS Frame 
     \item Server dekodiert HPACK-kodierte Header 
     \item Server aktualisiert die dynamische Dekodierungstabelle mit den neuen Headern
     \item Server verarbeitet die Anfrage und bereitet die Antwort mit HPACK-Kodierung mit der dynamischen Kodierungstabelle vor
\end{enumerate}

Antworten Verarbeiten

Der Client wartet bis alle Frames mit den entsprechenden Stream-IDs empfangen wurden und setzt die HTTP-Nachricht zusammen. Sobald die Nachricht vollständig ist, verarbeitet der Client die Antwort entsprechend der HTTP/2-Spezifikation, z.B. durch Anzeigen der Webseite oder Speichern der Daten.

Streams verwalten

Jede neue Anfrage öffnet einen neuen Stream mit einer neuen Stream-ID. Ein Thread könnte pro Stream laufen und den HTTP/2-Frames erstellen.
Dann wird auf der anderen Seite das Paket empfangen und basierend auf der Stream-ID zugeordnet und verarbeitet. Sobald die Antwort vollständig ist, wird der Stream geschlossen.
